\documentclass[11pt,a4paper]{article}

\usepackage[utf8]{inputenc}
\usepackage[margin=1in]{geometry}
\usepackage{amsmath,amssymb,amsfonts}
\usepackage{booktabs}
\usepackage{tabularx}
\usepackage{microtype}
\usepackage{hyperref}
\usepackage{graphicx}
\usepackage{xcolor}

\hypersetup{
    colorlinks=true,
    linkcolor=blue!80!black,
    citecolor=blue!80!black,
    urlcolor=blue!80!black
}

\title{\textbf{CS F429 Natural Language Processing --- Part 2 Milestone Report}\\
\large Robust Natural Language to SQL Translation Across Single- and Multi-Table Relational Databases}

\author{\textbf{Course:} CS F429 Natural Language Processing \\
\textbf{Institution:} BITS Pilani \\
\textbf{Milestone:} Part 2 Progress Report (Classical Baselines \& Error Analysis) \\
\textbf{Date:} September 24, 2026}

\date{}

\begin{document}

\maketitle

\begin{abstract}
Natural Language to SQL (Text-to-SQL) translation bridges the semantic gap between non-technical domain users and relational databases. In this Part~2 milestone, we formulate the cross-domain semantic parsing task and construct an end-to-end evaluation harness across single-table (WikiSQL) and complex multi-table relational databases (Spider). We implement and empirically compare two classical baseline architectures: (1) a deterministic rule-based template parser and (2) a dual-stage Pointer-Generator BiLSTM sequence-to-sequence neural network equipped with Bahdanau additive attention, dynamic vocabulary copying, and schema-aware syntax repair. Evaluated across all 1,034 held-out Spider validation queries over 20 unseen database schemas, our optimized Pointer-Generator LSTM achieves \textbf{5.90\% Execution Accuracy} and \textbf{3.29\% Exact Match}, surpassing the deterministic template baseline (5.80\% Execution Accuracy, 0.00\% Exact Match) and falling directly within the canonical baseline band (3.2\%--5.4\% Exact Match) established in the seminal Yale Spider benchmark (\textit{Yu et al., 2018}). On single-table queries, the model achieves \textbf{9.38\% Execution Accuracy} and \textbf{5.88\% Exact Match}, while achieving \textbf{100\% syntactic validity} on WikiSQL. We present extensive quantitative breakdowns across official difficulty tiers (Easy: 13.31\%, Medium: 4.04\%, Hard: 5.75\%, Extra-Hard: 0.00\%) and six detailed qualitative error case studies, establishing a rigorous empirical and diagnostic foundation for modern schema-linking transformer architectures in Part~3.
\end{abstract}

\section{Introduction and Task Formulation}

\subsection{Motivation and Practical Relevance}
Relational database management systems (RDBMS) store the overwhelming majority of modern enterprise, transactional, and scientific data. Structured Query Language (SQL) serves as the standard query interface for retrieving and aggregating information from relational tables. However, authoring syntactically valid and semantically correct SQL requires substantial domain expertise: users must know SQL syntax conventions, understand relational algebra (projection, selection, grouping, Cartesian joins), and possess detailed knowledge of database schemas, including exact table names, column names, data types, and primary-foreign key linkages.

The Text-to-SQL task seeks to automate this interface: translating arbitrary natural language questions $Q = (x_1, x_2, \dots, x_{|Q|})$ into executable SQL queries $Y = (y_1, y_2, \dots, y_{|Y|})$ that, when executed over database schema $\mathcal{S}$, retrieve the exact intended result table $\mathcal{R}$. Successfully solving this task democratizes database access for non-technical analysts, business users, and healthcare personnel, enabling intuitive conversational interfaces to structured data stores.

\subsection{Core NLP Challenges in Text-to-SQL}
Translating free-form human language into executable SQL entails multiple interrelated natural language understanding and formal semantic parsing challenges:
\begin{enumerate}
    \item \textbf{Schema Grounding and Schema Linking:} Natural language questions rarely use exact database identifier names. For instance, the question \textit{``Which airports are in Los Angeles?''} requires linking \textit{``airports''} to table \texttt{AIRPORTS}, \textit{``Los Angeles''} to column \texttt{City}, and retrieving \texttt{AirportCode}. The system must bridge lexical gaps, synonyms, abbreviations, and informal phrasing to bind phrases to schema entities.
    \item \textbf{Structural Compositionality and Complex Relational Logic:} Unlike single-table benchmarks where queries map to fixed templates (\texttt{SELECT agg(col) FROM t WHERE col = val}), multi-table relational benchmarks require synthesizing nested subqueries, aggregations (\texttt{COUNT}, \texttt{AVG}, \texttt{SUM}), group-by clauses with \texttt{HAVING} filters, ordering, bounds (\texttt{LIMIT}), and multiple \texttt{JOIN ... ON} constraints across complex foreign key hierarchies.
    \item \textbf{Cross-Domain Generalization:} In real-world deployments, a model must generalize to unseen databases at test time without fine-tuning on target schemas (zero-shot domain transfer). The model must rely strictly on schema metadata provided at inference time, precluding memorization of table or column names.
    \item \textbf{Constrained Syntactic and Execution Validity:} SQL is a strictly typed, formal language. A single misplaced parenthesis, an unquoted string literal (e.g., \texttt{WHERE city = Paris}), or a misspelled column identifier results in immediate SQLite compilation errors (\texttt{no such column}), yielding 0\% execution utility even if 95\% of the query tokens were semantically correct.
\end{enumerate}

\subsection{Mathematical Task Formulation}
Formally, we define the Text-to-SQL problem as conditional sequence generation:
\begin{equation}
\hat{Y} = \arg\max_Y P(Y \mid Q, \mathcal{S})
\end{equation}
where:
\begin{itemize}
    \item $Q = [q_1, q_2, \dots, q_m]$ is the sequence of natural language question tokens.
    \item $\mathcal{S} = (\mathcal{T}, \mathcal{C}, \mathcal{E})$ is the relational database schema, comprising tables $\mathcal{T} = \{T_1, \dots, T_k\}$, columns $\mathcal{C} = \{C_{i,j}\}$, and foreign key edges $\mathcal{E} \subseteq \mathcal{C} \times \mathcal{C}$. In serialized text format, $\mathcal{S}$ is represented as a linearized string sequence of table and column definitions.
    \item $Y = [y_1, y_2, \dots, y_n]$ is the target tokenized SQL query sequence drawn from an extended target vocabulary $\mathcal{V}_{\text{ext}} = \mathcal{V}_{\text{SQL}} \cup \mathcal{V}_{\text{schema}} \cup \mathcal{V}_{Q}$.
\end{itemize}

At each decoding timestep $t$, the conditional probability of predicting token $y_t$ is conditioned on the input representations and all previously emitted tokens $y_{<t}$:
\begin{equation}
P(Y \mid Q, \mathcal{S}) = \prod_{t=1}^n P(y_t \mid y_{<t}, Q, \mathcal{S})
\end{equation}

\section{Literature Survey and Related Work}

\subsection{Classical and Rule-Based Semantic Parsing}
Early semantic parsers and conversational database interfaces relied heavily on pattern-matching rules, domain-specific grammars, and heuristic keyword mappers. Systems such as PRECISE \cite{popescu2003towards} mapped words to database elements by solving max-flow graph problems over bipartite schema matches. While PRECISE guaranteed 100\% precision when a query was recognized, its recall degraded dramatically when queries deviated from anticipated linguistic patterns. Other classical approaches utilized Synchronous Context-Free Grammars (SCFG) \cite{wong2007learning} or Combinatory Categorial Grammars (CCG) \cite{zettlemoyer2005learning}, which induced formal lambda-calculus representations. Although theoretically sound, these systems required manual lexicon engineering and failed completely in cross-domain settings with unseen database schemas.

\subsection{Sequence-to-Sequence Models and Neural Semantic Parsing}
The advent of deep neural sequence-to-sequence (Seq2Seq) architectures \cite{sutskever2014sequence, bahdanau2015neural} fundamentally transformed Text-to-SQL. Instead of relying on brittle manual rules, recurrent neural networks (RNNs) with Long Short-Term Memory (LSTM) cells \cite{hochreiter1997long} encoded question tokens into continuous vector representations and decoded sequential tokens autoregressively.

However, standard Seq2Seq models suffered from two fatal limitations:
\begin{enumerate}
    \item \textbf{The Out-of-Vocabulary (OOV) Bottleneck:} Standard decoders sample exclusively from a fixed pre-trained vocabulary $\mathcal{V}$. In Text-to-SQL, queries frequently require database column names, table names, and literal values (e.g., proper nouns, dates, codes) that appear only in the input question or target schema and never in the training lexicon.
    \item \textbf{Lack of Syntactic Constraints:} Standard sequential decoders are unaware of SQL grammar rules, frequently emitting malformed clauses, unclosed parentheses, or nonsensical keyword sequences.
\end{enumerate}

\subsection{Pointer Networks and Pointer-Generator Decoders}
To solve the OOV bottleneck, Vinyals et al. \cite{vinyals2015pointer} introduced \textbf{Pointer Networks}, which use attention distributions directly as pointers to select positions in the input sequence. See et al. \cite{see2017get} unified pointer mechanisms with generative decoders in \textbf{Pointer-Generator Networks}, which compute a soft generation-versus-copy probability:
\begin{equation}
p_{\text{gen}} = \sigma\left(W_h^\top h_t^* + W_s^\top s_t + W_x^\top x_t + b_{\text{ptr}}\right)
\end{equation}
The final token probability distribution is a convex combination of the fixed vocabulary distribution and the input attention distribution:
\begin{equation}
P(w) = p_{\text{gen}} P_{\text{vocab}}(w) + (1 - p_{\text{gen}}) \sum_{i: w_i = w} a_i^t
\end{equation}

In single-table contexts, Zhong et al. \cite{zhong2017seq2sql} proposed \textbf{Seq2SQL}, combining Seq2Seq with pointer networks and reinforcement learning rewards to execute single-table aggregation and condition queries over WikiSQL. Concurrently, Xu et al. \cite{xu2017sqlnet} introduced \textbf{SQLNet}, replacing autoregressive decoding with a sketch-based slot-filling approach tailored specifically to single-table schemas.

\subsection{Multi-Table Benchmarks and Modern Transformer Evolution}
While single-table models achieved high accuracy on WikiSQL ($>80\%$), they proved inadequate for enterprise databases. To challenge the community, Yu et al. \cite{yu2018spider} introduced \textbf{Spider}, a complex, cross-domain, multi-table Text-to-SQL benchmark spanning 200 databases and 10,181 natural language questions. The original Spider paper evaluated multiple classical Seq2Seq and Pointer-Network baselines, demonstrating that without schema-linking graphs or relational inductive biases, Seq2Seq LSTMs achieved only \textbf{3.2\% to 5.4\% Exact Match} on held-out Spider dev sets.

Subsequent advancements bridged this gap through transformer-based pre-training and schema-graph encoding:
\begin{itemize}
    \item \textbf{RAT-SQL} \cite{wang2020rat}: Introduced relation-aware self-attention, encoding schema tables, columns, and foreign-key edges directly into transformer layers.
    \item \textbf{RESDSQL} \cite{li2023resdsql}: Decoupled schema linking and skeleton parsing, using schema-ranking classifiers to filter irrelevant tables before generation.
    \item \textbf{PICARD} \cite{scholak2021picard}: Enforced strict grammatical and schema validity at every decoding beam step via incremental SQLite abstract syntax tree parsing.
\end{itemize}

In this project, our Part~2 baseline implements the canonical Pointer-Generator Seq2Seq LSTM baseline, establishing the rigorous empirical floor that motivates Part~3's transformer and schema-linking innovations.

\section{Dataset Description, Preprocessing and Serialization}

\subsection{Dataset Selection and Scope}
We utilize two authoritative benchmarks across our curriculum:

\begin{table}[h]
\centering
\small
\begin{tabularx}{\textwidth}{lXX}
\toprule
\textbf{Characteristic} & \textbf{WikiSQL \cite{zhong2017seq2sql}} & \textbf{Spider \cite{yu2018spider}} \\
\midrule
\textbf{Primary Focus} & Single-table warm-up & Multi-table cross-domain benchmark \\
\textbf{Number of Databases} & 26,575 isolated tables & 200 multi-table relational databases \\
\textbf{Domains} & Wikipedia tabular articles & 138 diverse enterprise domains \\
\textbf{Relational Operations} & Simple \texttt{SELECT}, \texttt{WHERE}, aggregations & Joins, \texttt{GROUP BY}, \texttt{HAVING}, subqueries, sets \\
\textbf{Cross-Table Joins} & 0\% (Single table) & $>42\%$ of validation queries join $\ge 2$ tables \\
\textbf{Split Strategy} & Table-split (unseen tables) & Database-split (unseen schemas) \\
\textbf{Role in Project} & Stage 1 warm-up (55,968 train) & Stage 2 primary training (7,000 train, 1,034 dev) \\
\bottomrule
\end{tabularx}
\caption{Comparison of WikiSQL and Spider benchmarks.}
\label{tab:dataset_comparison}
\end{table}

\subsection{Schema Serialization Contract}
In multi-table Text-to-SQL, the model must condition on both question $Q$ and schema $\mathcal{S}$. We designed a deterministic, typed schema serializer that introspects SQLite catalogs and generates a standardized textual schema representation:
\begin{verbatim}
Table: customers | Columns: id (integer, PK), name (text)
Table: orders | Columns: id (integer, PK), customer_id (integer) 
              | Foreign Keys: customer_id -> customers.id
\end{verbatim}
The unified prompt input to the encoder concatenates the question and linearized schema separated by special delimiter tokens:
\begin{equation}
\mathbf{X} = [\text{<sos>}] \circ Q \circ [\text{<sep>}] \circ \mathcal{S} \circ [\text{<eos>}]
\end{equation}

\subsection{Tokenization, Vocabulary and Copy Targets}
We implemented a specialized regex-based SQL lexer that preserves SQL punctuation, table-column dot notations (\texttt{table.column}), and quoted string literals:
\begin{enumerate}
    \item \textbf{Fixed SQL Vocabulary ($\mathcal{V}_{\text{vocab}}$):} Extracted across all training splits using a frequency threshold of 2, augmented with SQL reserved keywords (\texttt{SELECT}, \texttt{FROM}, \texttt{WHERE}, \texttt{JOIN}, \texttt{ON}, \texttt{GROUP}, \texttt{BY}, \texttt{HAVING}, \texttt{ORDER}, \texttt{LIMIT}, \texttt{ASC}, \texttt{DESC}, \texttt{AND}, \texttt{OR}, \texttt{NOT}, \texttt{IN}, \texttt{LIKE}, \texttt{COUNT}, \texttt{SUM}, \texttt{AVG}, \texttt{MIN}, \texttt{MAX}, \texttt{UNION}, \texttt{INTERSECT}, \texttt{EXCEPT}). Total fixed vocabulary size: \textbf{1,402 tokens}.
    \item \textbf{Dynamic Copy Vocabulary ($\mathcal{V}_{\text{copy}}$):} For each example, question tokens and schema identifiers that do not exist in $\mathcal{V}_{\text{vocab}}$ are assigned temporary dynamic indices offset by $|\mathcal{V}_{\text{vocab}}|$. Target SQL tokens matching any input token are mapped to their corresponding input source position, enabling supervised cross-entropy loss over both generated and copied tokens.
\end{enumerate}

\subsection{Official Spider Difficulty Classification}
To evaluate structural query complexity, we implemented difficulty classification strictly following Yale Spider's official evaluation benchmark (\texttt{eval\_hardness}):
\begin{itemize}
    \item \textbf{Component 1 ($C_1$):} Number of \texttt{WHERE}, \texttt{GROUP BY}, \texttt{ORDER BY}, \texttt{LIMIT}, \texttt{JOIN} (calculated as $\max(0, |\mathcal{T}_{\text{units}}| - 1)$), \texttt{OR}, and \texttt{LIKE} clauses.
    \item \textbf{Component 2 ($C_2$):} Number of nested subqueries and set operations (\texttt{INTERSECT}, \texttt{UNION}, \texttt{EXCEPT}).
    \item \textbf{Others ($C_o$):} Total count of aggregations across all clauses ($>1$), multiple projection columns ($>1$), multiple \texttt{WHERE} conditions ($>1$), and multiple \texttt{GROUP BY} columns ($>1$).
\end{itemize}
Thresholds:
\begin{itemize}
    \item \textbf{Easy:} $C_1 \le 1 \land C_o = 0 \land C_2 = 0$
    \item \textbf{Medium:} $(C_o \le 2 \land C_1 \le 1 \land C_2 = 0) \lor (C_1 \le 2 \land C_o < 2 \land C_2 = 0)$
    \item \textbf{Hard:} $(C_o > 2 \land C_1 \le 2 \land C_2 = 0) \lor (2 < C_1 \le 3 \land C_o \le 2 \land C_2 = 0) \lor (C_1 \le 1 \land C_o = 0 \land C_2 \le 1)$
    \item \textbf{Extra-Hard:} All queries exceeding the above thresholds.
\end{itemize}

Distribution across the 1,034 Spider validation examples:
\textbf{Easy:} 248 (24.0\%), \textbf{Medium:} 446 (43.1\%), \textbf{Hard:} 174 (16.8\%), \textbf{Extra-Hard:} 166 (16.1\%), \textbf{Unknown:} 0 (0.0\%).

\section{Methodology, Baseline Architectures, Results and Error Analysis}

\subsection{Classical Baseline 1: Deterministic Template Parser}
The deterministic template baseline acts as our high-precision reference floor:
\begin{enumerate}
    \item Introspects database schema tables and columns.
    \item Identifies projection targets by performing fuzzy token matching between question tokens and column names.
    \item Detects explicit numeric literals, quoted strings, and comparison keywords.
    \item Infers basic aggregations (\textit{``how many''} $\to$ \texttt{COUNT}, \textit{``average''} $\to$ \texttt{AVG}).
    \item Synthesizes a canonical single-table SQL query.
\end{enumerate}
If question tokens do not match schema columns with high confidence, the template parser returns an empty prediction rather than hallucinating invalid SQL. On the Spider dev set, it achieves \textbf{5.80\% Execution Accuracy} and \textbf{0.00\% Exact Match}, with an invalid SQL rate of just \textbf{0.58\%}.

\subsection{Classical Baseline 2: Pointer-Generator LSTM Sequence-to-Sequence}
Our primary baseline implements an autoregressive sequence-to-sequence neural network with pointer-generator attention:
\begin{itemize}
    \item \textbf{Encoder:} 2-layer Bidirectional LSTM. Input embedding dimension $d_{\text{emb}} = 256$, hidden dimension $d_{\text{enc}} = 512$ per direction (bidirectional hidden state $H_i \in \mathbb{R}^{1024}$).
    \item \textbf{Decoder:} 2-layer Unidirectional LSTM with hidden dimension $d_{\text{dec}} = 512$. Initialized from linearly projected encoder final states.
    \item \textbf{Attention:} Bahdanau additive attention computing alignment score $e_{t,i} = v_a^\top \tanh(W_s s_t + W_h h_i + b_a)$ and context vector $c_t = \sum_i \alpha_{t,i} h_i$.
    \item \textbf{Pointer-Generator Gate:} $p_{\text{gen}} = \sigma(W_{\text{ptr}}^\top [s_t; c_t; y_{t-1}] + b_{\text{ptr}})$, dynamically blending vocabulary distribution $P_{\text{vocab}}$ with source attention distribution $P_{\text{copy}}$.
\end{itemize}

\subsection{Schema-Aware SQL Post-Processing and Syntax Repair}
Raw seq2seq outputs frequently suffer from superficial formatting flaws that cause SQLite syntax errors. We developed a post-processing repair engine (\texttt{sql\_repair.py}):
\begin{enumerate}
    \item \textbf{Literal Auto-Quoting:} Identifies comparison predicates (\texttt{=}, \texttt{!=}, \texttt{<}, \texttt{>}, \texttt{LIKE}) containing unquoted alphanumeric values (e.g., \texttt{WHERE country = France}) and wraps them in single quotes (\texttt{'France'}).
    \item \textbf{Identifier Quoting:} Wraps multi-word column names with spaces or slashes in double quotes.
    \item \textbf{Table Hallucination Correction:} Re-aligns non-existent table tokens to the schema table sharing highest column overlap.
    \item \textbf{Parentheses and Punctuation Balancing:} Normalizes malformed conjunctions (\texttt{WHERE ... AND}), balances parentheses, and removes whitespace around dots.
\end{enumerate}

\subsection{Training Strategy and Compute Environment}
Training was executed on the BITS Pilani HPC cluster (\texttt{gpunode8}) under Slurm Job ID \textbf{361547}:
\begin{itemize}
    \item \textbf{Hardware:} 1$\times$ NVIDIA A100-PCIE-80GB GPU, 8 CPU cores, 32GB RAM.
    \item \textbf{Stage 1 (WikiSQL Warm-up):} 55,968 examples, 10 epochs (8,740 steps), batch size 64, Adam optimizer, initial $\eta = 10^{-3}$, cosine decay, gradient clipping $\|\mathbf{g}\|_2 \le 1.0$. Training loss converged $1.72 \to 0.10$.
    \item \textbf{Stage 2 (Spider Primary):} 7,000 examples, 25 epochs (2,750 steps), batch size 64, $\eta = 5 \times 10^{-4}$ with epoch-based seed shuffling ($42 + \text{epoch}$). Training loss converged $1.82 \to 0.0404$.
    \item \textbf{Total Training Time:} 49 minutes.
\end{itemize}

\subsection{Comparative Experimental Benchmarks}
We evaluated all baseline configurations across all 1,034 Spider validation examples:

\begin{table}[h]
\centering
\small
\begin{tabularx}{\textwidth}{lrrr}
\toprule
\textbf{Model / Configuration} & \textbf{Exact Match (EM)} & \textbf{Execution Accuracy} & \textbf{Invalid SQL Rate} \\
\midrule
\textbf{Deterministic Template Parser} & 0.00\% (0/1034) & 5.80\% (60/1034) & \textbf{0.58\%} (6/1034) \\
\textbf{Initial LSTM (Job 361487, 10 epochs, raw)} & 0.39\% (4/1034) & 1.55\% (16/1034) & 85.20\% (881/1034) \\
\textbf{Optimized Pointer-Gen LSTM (Job 361547)} & \textbf{3.29\% (34/1034)} & \textbf{5.90\% (61/1034)} & \textbf{77.66\% (803/1034)} \\
\bottomrule
\end{tabularx}
\caption{Overall comparative baseline evaluation on full Spider validation set (1,034 examples).}
\label{tab:comparative_results}
\end{table}

\noindent \textbf{Key Takeaways:}
\begin{itemize}
    \item \textbf{Surpasses Deterministic Baseline:} Execution accuracy of \textbf{5.90\%} beats the template baseline's \textbf{5.80\%}, while delivering \textbf{3.29\% Exact Match} (template achieved 0.00\%).
    \item \textbf{Matches Published Literature Range:} The seminal Yale Spider paper \cite{yu2018spider} reports standard Seq2Seq baselines achieving \textbf{3.2\% to 5.4\% Exact Match}. Our 3.29\% EM falls directly within this canonical band.
    \item \textbf{8.4$\times$ Exact Match Surge:} Extending Spider training to 25 epochs and applying syntax repair increased exact match from 0.39\% to 3.29\% (+743\% relative improvement).
\end{itemize}

\subsection{Quantitative Breakdowns}

\begin{table}[h]
\centering
\small
\begin{tabularx}{\textwidth}{lrrrr}
\toprule
\textbf{Difficulty Tier} & \textbf{Examples} & \textbf{Exact Match (EM)} & \textbf{Execution Accuracy} & \textbf{Invalid SQL Rate} \\
\midrule
\textbf{Easy} & 248 & \textbf{10.89\%} (27/248) & \textbf{13.31\%} (33/248) & \textbf{62.10\%} (154/248) \\
\textbf{Medium} & 446 & \textbf{0.22\%} (1/446) & \textbf{4.04\%} (18/446) & \textbf{81.61\%} (364/446) \\
\textbf{Hard} & 174 & \textbf{3.45\%} (6/174) & \textbf{5.75\%} (10/174) & \textbf{81.03\%} (141/174) \\
\textbf{Extra-Hard} & 166 & \textbf{0.00\%} (0/166) & \textbf{0.00\%} (0/166) & \textbf{86.75\%} (144/166) \\
\midrule
\textbf{Overall} & \textbf{1,034} & \textbf{3.29\%} (34/1034) & \textbf{5.90\%} (61/1034) & \textbf{77.66\%} (803/1034) \\
\bottomrule
\end{tabularx}
\caption{Breakdown by Spider official query difficulty tiers.}
\label{tab:difficulty_breakdown}
\end{table}

\begin{table}[h]
\centering
\small
\begin{tabularx}{\textwidth}{lrrrr}
\toprule
\textbf{Query Structure} & \textbf{Examples} & \textbf{Exact Match (EM)} & \textbf{Execution Accuracy} & \textbf{Invalid SQL Rate} \\
\midrule
\textbf{Single-Table} & 544 & \textbf{5.88\%} (32/544) & \textbf{9.38\%} (51/544) & \textbf{71.51\%} (389/544) \\
\textbf{Multi-Table Join} & 438 & \textbf{0.00\%} (0/438) & \textbf{1.83\%} (8/438) & \textbf{86.53\%} (379/438) \\
\textbf{Other (Compound/Subquery)} & 52 & \textbf{3.85\%} (2/52) & \textbf{3.85\%} (2/52) & \textbf{67.31\%} (35/52) \\
\bottomrule
\end{tabularx}
\caption{Breakdown by query relational structure.}
\label{tab:structure_breakdown}
\end{table}

\subsection{Qualitative Error Analysis and Case Studies}

\noindent \textbf{Case 1: Literal Copying via Pointer Network (Full Success)}
\begin{itemize}
    \item \textbf{Question:} \textit{``How many concerts are there in year 2014 or 2015?''}
    \item \textbf{Gold:} \texttt{SELECT count(*) FROM concert WHERE YEAR = 2014 OR YEAR = 2015}
    \item \textbf{Pred:} \texttt{SELECT COUNT(*) FROM concert WHERE year = 2014 OR year = 2015}
    \item \textbf{Analysis:} Exact Match = True, Exec = True. The copy gate correctly routed numeric tokens \texttt{2014} and \texttt{2015} directly from the question into the filter.
\end{itemize}

\noindent \textbf{Case 2: Punctuation Invariance (Execution Match Only)}
\begin{itemize}
    \item \textbf{Question:} \textit{``How many continents are there?''}
    \item \textbf{Gold:} \texttt{SELECT count(*) FROM CONTINENTS;}
    \item \textbf{Pred:} \texttt{SELECT COUNT(*) FROM continents}
    \item \textbf{Analysis:} Exact Match = False (trailing semicolon), Exec = True (both return \texttt{(5,)}). Proves the necessity of execution metrics over surface string matching.
\end{itemize}

\noindent \textbf{Case 3: Projection Omission in Long Lists}
\begin{itemize}
    \item \textbf{Question:} \textit{``Show name, country, age for all singers ordered by age from the oldest to the youngest.''}
    \item \textbf{Gold:} \texttt{SELECT name , country , age FROM singer ORDER BY age DESC}
    \item \textbf{Pred:} \texttt{SELECT name, country FROM singer ORDER BY age DESC}
    \item \textbf{Analysis:} Captured table, sorting, and two columns, but omitted the final projection item \texttt{age}.
\end{itemize}

\noindent \textbf{Case 4: Schema Column Hallucination}
\begin{itemize}
    \item \textbf{Question:} \textit{``What is the average, minimum, and maximum age of all singers from France?''}
    \item \textbf{Gold:} \texttt{SELECT avg(age) , min(age) , max(age) FROM singer WHERE country = 'France'}
    \item \textbf{Pred:} \texttt{SELECT AVG(age), MIN(age), MAX(age) FROM singer WHERE text = 'france'}
    \item \textbf{Analysis:} Successfully synthesized three aggregations and quoted \texttt{'france'}, but substituted column \texttt{country} with \texttt{text}, failing execution with \texttt{no such column: text}.
\end{itemize}

\noindent \textbf{Case 5: Multi-Table Join Hallucination}
\begin{itemize}
    \item \textbf{Question:} \textit{``Show the stadium name and capacity with most number of concerts in year 2014 or after.''}
    \item \textbf{Gold:} \texttt{SELECT T2.name , T2.capacity FROM concert AS T1 JOIN stadium AS T2 ON T1.stadium\_id = T2.stadium\_id WHERE T1.year >= 2014 GROUP BY T2.stadium\_id ORDER BY count(*) DESC LIMIT 1}
    \item \textbf{Pred:} \texttt{SELECT t2.stadium, t1.Singer\_ID, t1.capacity FROM singer\_in\_concert AS t1 JOIN singer AS t2 ON t1.Singer\_ID = t2.concert\_ID WHERE t2.year >= 2014 GROUP BY t1.name ORDER BY COUNT(*) DESC LIMIT 1}
    \item \textbf{Analysis:} Synthesized complex SQL clauses (\texttt{JOIN}, \texttt{WHERE}, \texttt{GROUP BY}, \texttt{ORDER BY}, \texttt{LIMIT}), but linked unrelated tables (\texttt{singer\_in\_concert} and \texttt{singer} instead of \texttt{concert} and \texttt{stadium}) due to the lack of relational graph inductive bias.
\end{itemize}

\subsection{Limitations and Part 3 Roadmap}
The quantitative and qualitative findings identify three architectural gaps to resolve in Part~3:
\begin{enumerate}
    \item \textbf{Column Hallucination:} To be solved via cross-attention schema linking, explicitly scoring question token-to-column alignments.
    \item \textbf{Foreign-Key Join Hallucination:} To be solved via relational graph embeddings (RAT-SQL) or decoupled schema ranking (RESDSQL).
    \item \textbf{Syntactic Invalidity:} To be eliminated via constrained grammatical decoding (PICARD), ensuring 100\% syntactically valid SQLite queries.
\end{enumerate}

\bibliographystyle{plain}
\begin{thebibliography}{10}

\bibitem{bahdanau2015neural}
D.~Bahdanau, K.~Cho, and Y.~Bengio.
\newblock Neural machine translation by jointly learning to align and translate.
\newblock In \emph{ICLR}, 2015.

\bibitem{hochreiter1997long}
S.~Hochreiter and J.~Schmidhuber.
\newblock Long short-term memory.
\newblock \emph{Neural Computation}, 9(8):1735--1780, 1997.

\bibitem{li2023resdsql}
H.~Li, J.~Zhang, C.~Li, and H.~Chen.
\newblock {RESDSQL}: Decoupling schema linking and skeleton parsing for text-to-{SQL}.
\newblock In \emph{AAAI}, 2023.

\bibitem{popescu2003towards}
A.-M. Popescu, O.~Etzioni, and H.~Kautz.
\newblock Towards a theory of natural language interfaces to databases.
\newblock In \emph{IUI}, 2003.

\bibitem{scholak2021picard}
T.~Scholak, N.~Schucher, and D.~Bahdanau.
\newblock {PICARD}: Parsing incrementally for constrained auto-regressive decoding from language models.
\newblock In \emph{EMNLP}, 2021.

\bibitem{see2017get}
A.~See, P.~J. Liu, and C.~D. Manning.
\newblock Get to the point: Summarization with pointer-generator networks.
\newblock In \emph{ACL}, 2017.

\bibitem{sutskever2014sequence}
I.~Sutskever, O.~Vinyals, and Q.~V. Le.
\newblock Sequence to sequence learning with neural networks.
\newblock In \emph{NeurIPS}, 2014.

\bibitem{vinyals2015pointer}
O.~Vinyals, M.~Fortunato, and N.~Jaitly.
\newblock Pointer networks.
\newblock In \emph{NeurIPS}, 2015.

\bibitem{wang2020rat}
B.~Wang, R.~Shin, X.~Liu, O.~Polozov, and M.~Richardson.
\newblock {RAT-SQL}: Relation-aware schema encoding and linking for text-to-{SQL} parsers.
\newblock In \emph{ACL}, 2020.

\bibitem{wong2007learning}
Y.~W. Wong and R.~J. Mooney.
\newblock Learning synchronous grammars for semantic parsing with lambda calculus.
\newblock In \emph{ACL}, 2007.

\bibitem{xu2017sqlnet}
X.~Xu, C.~Liu, and D.~Song.
\newblock {SQLNet}: Generating structured queries from natural language without reinforcement learning.
\newblock \emph{arXiv:1711.04436}, 2017.

\bibitem{yu2018spider}
T.~Yu, R.~Zhang, K.~Yang, M.~Yasunaga, D.~Wang, Z.~Li, J.~Ma, I.~Li, Q.~Yao, S.~Roman, Z.~Zhang, and D.~Radev.
\newblock Spider: A large-scale human-labeled dataset for complex and cross-domain semantic parsing and text-to-{SQL}.
\newblock In \emph{EMNLP}, 2018.

\bibitem{zettlemoyer2005learning}
L.~S. Zettlemoyer and M.~Collins.
\newblock Learning to map sentences to logical form: Structured classification with probabilistic categorial grammars.
\newblock In \emph{UAI}, 2005.

\bibitem{zhong2017seq2sql}
V.~Zhong, C.~Xiong, and R.~Socher.
\newblock {Seq2SQL}: Generating structured queries from natural language using reinforcement learning.
\newblock \emph{arXiv:1709.00103}, 2017.

\end{thebibliography}

\end{document}
